\chapter{Sobre este Livro}

Este livro é um Recurso Educacional Aberto. Ele foi escrito para estudantes
com algum conhecimento de lógica e de matemática do ensino médio.
Seu objetivo é oferecer uma primeira visão da filosofia da teoria dos
conjuntos. Ao final deste livro, os estudantes que o lerem podem ter uma
noção:
\begin{enumerate}
	\item de por que a teoria dos conjuntos surgiu;
	\item de como incorporar grandes partes da matemática em teoria dos conjuntos + aritmética;
	\item de como incorporar a própria aritmética na teoria dos conjuntos;
	\item do que envolve a concepção cumulativa iterativa de conjunto;
	\item de como se pode tentar justificar os axiomas de ZFC.
\end{enumerate}
O livro surgiu de um curso curto que eu ministrei no Departamento de Filosofia de Cambridge. Antes de mim, ele foi ministrado por Luca Incurvati e Michael
Potter. Ao escrever este livro---e, de forma mais geral, o curso---fiquei
profundamente em dívida com Luca e Michael. Espero que isso fique claro no próprio texto do livro; mas também quero registrar aqui
meus sinceros agradecimentos a ambos.

A maior parte deste livro foi originalmente publicada como \emph{Open
Set Theory}; este livro é seu sucessor. Contribuí com todo o material deste livro para o \href{http://openlogicproject.org}{\emph{Open Logic Project}}, onde ele está disponível gratuitamente. (\Crefrange{sfr:set::chap}{sfr:siz::chap} foram elaborados, com apenas
pequenas alterações, a partir de material já existente no OLP; essas alterações agora também fazem parte do OLP.) Veja
\href{http://openlogicproject.org/}{openlogicproject.org} para mais
informações.

\begin{flushright}
Tim Button
\\University College London
\\Outubro de 2021
\end{flushright}
